% % % % % % % % % % % % % % % % % % % % % % % % % % % % % % % % % % % % % % % % % % 
% Exemple de rapport de master
% Version originale :  A. Tarantola, Octobre 2004
% Version modifiée :   A. Fournier et G. Moguilny, Mars 2012... Novembre 2025
% % % % % % % % % % % % % % % % % % % % % % % % % % % % % % % % % % % % % % % % % % 

\documentclass{ipgpmaster}
%
% Pour saisir directement les lettres accentuées :
% \usepackage[applemac]{inputenc} % pour utilisateurs mac
% \usepackage[ansinew]{inputenc}  % pour Windows ANSI
% \usepackage[latin1]{inputenc}   % pour Unix Latin1
\usepackage[T1]{fontenc}
\usepackage{amsmath}

% \bibliography{master_AT}  % pour bibtex
\addbibresource{masterBib.bib} % pour biblatex

\begin{document}

\checkyears

\vspace*{5mm}

%\setkeys{Gin}{draft} % images non affichées
\setkeys{Gin}{draft=false} % images affichées

%\linenumbers % Numerotation des lignes pour relecture
%Select here the language 
\selectlanguage{french}   
%\selectlanguage{english} 


\def\author{Albert Tarantola}
%\footnote{
%Albert Tarantola (1949-2009) : Géophysicien français d'origine espagnole, dont les travaux %ont notamment porté sur la théorie du problème
%inverse. Il a effectué sa carrière à l'IPGP, étant notamment responsable du DEA (diplôme %d'études approfondies, l'ancêtre du master 2) de géophysique interne entre
%1991 et 1996. En tant que responsable du DEA, il a encouragé l'utilisation de \LaTeX\ pour %la rédaction des rapports de stage, et rédigé les recommandations reprises dans ce document.
%}}
\def\title{Le Rapport de Stage de Master}
\def\shorttitle{Rapport de stage}
\def\unit{IPGP}
\def\team{Dynamique des fluides géologiques \& Sismologie}
\def\spe{Géophysique ou autre} % Parcours
\def\supervisor{A. Fournier \& G. Moguilny}
\def\mydate{\today}
\def\keywords{Sciences de la Terre,  master, rapport, conseils}

\Entete


\begin{abstract}
Le stage de recherche du master est essentiel dans la formation d'un scientifique, 
car il permet un premier contact avec le monde de la recherche. 
La rédaction du rapport est une composante importante du stage.
Il doit être rédigé en suivant le style littéraire habituel des articles 
scientifiques, avec ses usages bien codifiés. 
Les différentes parties du rapport : 
\emph{Résumé, Introduction, Résultats, Discussion, Conclusion, Remerciements, 
Références bibliographiques,\/} et \emph{Annexes,\/} 
ont chacune des objectifs bien dé\-ter\-mi\-nés. 
On doit chercher la perfection dans la forme.
\end{abstract}

\vfill

\rule{\linewidth}{0.4pt}\\[3mm]
\textit{
Albert Tarantola (1949-2009) : Géophysicien français d'origine espagnole, dont les travaux ont notamment porté sur la théorie du problème
inverse. Il a effectué sa carrière à l'IPGP, étant notamment responsable du DEA (diplôme d'études approfondies, l'ancêtre du master 2) de géophysique interne entre
1991 et 1996. En tant que responsable du DEA, il a encouragé l'utilisation de \LaTeX\ pour la rédaction des
rapports de stage, et rédigé des recommandations reprises dans ce document.
}


\newpage
\tableofcontents
\newpage

\section{Introduction}

Le stage de recherche du master permet aux étudiants d'apprendre le métier 
de chercheur, en participant au travail de recherche d'un laboratoire. 
Le travail d'un chercheur ressemble beaucoup à celui d'un journaliste : 
il faut d'abord avoir l'idée d'un sujet d'enquête, ensuite, il faut que 
l'enquête fournisse des informations et, enfin, il faut la publier. 
Une grande quantité d'enquêtes journalistiques, et de recherches 
scientifiques, sont restées ignorées parce que leur publication a été 
bâclée.

Rédiger un rapport scientifique prend beaucoup de temps. Après un 
``premier jet'', il faut souvent tout réécrire. 
Lorsque la forme du rapport semble acceptable, il faut demander l'avis 
de ses collègues, y apporter des modifications (souvent dramatiques), 
et itérer le processus
jusqu'à ce qu'un document de lecture facile ait été produit. Il~existe 
un \emph{style littéraire\/} très strict que les jeunes chercheurs ont parfois 
du mal à accepter. Si l'originalité est indispensable pour le sujet de 
recherche, elle est simplement distrayante dans la publication,
où on ne demande que de la perfection dans la pré\-sen\-ta\-tion.


\section{Les différentes parties d'un article}


\subsection{Le résumé}

Partie très importante du rapport.
Il doit permettre de se faire une idée précise des résultats.
Le style est identique à celui de l'article, simplement en résumé.
Oubliez que le résumé est publié au début de l'article et rédigez 
comme s'il devait être la seule chose qui existe : ne faites pas 
référence à l'article (si vous disiez ``nous montrons que\dots'' 
vous seriez en train de le faire). 
Ne parlez pas de ``certaines'' hypothèses ou résultats ; dites lesquels 
ou lesquelles (ou n'en parlez pas si ce n'est pas important). Il est difficile de véhiculer beaucoup 
d'information en peu de lignes, mais c'est indispensable.


\subsection{L'introduction}

Sert à situer le contexte du travail. Pourquoi est-ce intéressant ? 
Qui a travaillé avant sur ce sujet? Quelle démarche avons-nous suivies ?
Il faut aussi dans l'introduction expliquer l'organisation de l'article 
(après quelques rappels de la théorie de base,  l'instrument 
utilisé est décrit dans la section 3.2\dots), sans que cela 
devienne une ``table des matières''.


\subsection{La méthodologie, l'analyse, la discussion}

Évidemment, la partie la plus importante dans le fond, très difficile 
à écrire. C'est ici que la démarche scientifique doit être 
scrupuleusement suivie. De la rigueur, de la rigueur, et de la clarté.


\subsection{Les conclusions}

Les conclusions doivent déjà apparaître dans le résumé et le 
corps de l'article. Ici, en utilisant un langage technique ---permis par la
lecture de l'ensemble du rapport,--- on doit oublier les détails, faire ressortir l'essentiel et \emph{ouvrir des perspectives\/}.


\subsection{Les références bibliographiques}

Il existe plusieurs styles pour les références bibliographiques. 
Il en existe même trop : raison de plus pour ne pas en inventer davantage. 
Pour le rapport de stage du master, suivez le style 
défini avec ce gabarit. 

Notez que l'article de \cite{cao1990simultaneous} contient, dans 
la liste de références bi\-blio\-gra\-phiques, le nom explicite de tous 
les auteurs.
L'article de \cite{crase1992nonlinear} est sorti.  
 Il y a aussi un article qui n'a été que 
soumis \parencite{frezatsoumis}. 
En fait, si 
vous utilisez \LaTeX\ et ce gabarit, le formatage des références est géré 
automatiquement et vous n'avez pas à vous occuper du style 
(quand vous aurez besoin d'aide, adressez-vous à Geneviève Moguilny,
au troisième étage). 
%Pour plus de précisions sur \textsc{Bib}\LaTeX{} avec \texttt{natbib}, 
% consultez également la partie~\ref{sec:natbib}. 

Lors de votre présentation orale, souvenez-vous que la locution latine 
``et al.'' est réservée à l'écrit : on ne dit pas ``Cao et al.'', 
mais ``Cao et collaborateurs''.

\subsection{Les annexes}

Un article doit être très lisible. Pour cela, il est utile de mettre en 
annexe toute démonstration ou discussion un peu lourde qui peut 
être distrayante. Il est préférable de n'écrire dans l'article, 
dans un langage simple, que la trame des conclusions, qui seront justifiées 
dans le détail seulement dans les annexes.


\section{Quelques remarques d'ordre général}

L'information essentielle du rapport doit se trouver à {\bfseries quatre} endroits
différents : dans le résumé, dans le corps de l'article, dans les 
légendes des figures et des tables, et dans la conclusion.

Il faut lire beaucoup d'articles scientifiques. C'est la seule façon 
de saisir le \emph{style littéraire\/} d'un rapport scientifique.
N'envisagez pas d'écrire votre rapport sans une bonne connaissance des 
revues scientifiques importantes (\emph{Nature, Science, Journal of 
Geo\-phy\-si\-cal Research, Geophysical Journal International, 
Journal of Fluid Mechanics, Seismica, Journal of Studies of the Earth's Deep Interior, 
\dots\/}). Le style d'un rapport scientifique n'est pas celui des 
revues de vulgarisation (\emph{Epsiloon, La Recherche\dots\/})
qui, par ailleurs, peuvent être excellentes.

De plus, pour ce rapport, au format A4 :
\begin{itemize}
\item la première page doit contenir le résumé de 20 à 25 lignes, 
\item la deuxième page doit contenir le Sommaire,
\item le nombre de pages ne doit pas excéder 30,
\item les pages doivent être  numérotées en haut à droite/gauche (pour celles
de numéro impair/pair). 
\end{itemize}

\subsection{Les figures}

Si possible, les points clefs de l'article doivent être illustrés par une 
figure (celle qui vaut plus que 100~mots).

En utilisant des logiciels habituels de dessin (e.g., {\sl Illustrator\/}, 
 {\sl gimp\/}), 
il est possible de produire des figures propres qui peuvent s'intégrer 
dans l'article en utilisant un des
paquets (ou \textit{packages}) existants 
comme \texttt{graphicx}.
Essayez, dans la mesure du possible, de produire des figures en format 
\texttt{pdf} (\textit{Portable Document Format\/}).

Il n'est pas interdit d'utiliser une figure extraite d'un livre ou d'un 
article, mais il ne faut pas oublier d'en donner 
la référence exacte.


Dans la forme finale du rapport, les figures doivent être à leur place.
Les légendes des figures doivent être assez explicites pour permettre 
de comprendre le travail, même si on ne lit que le résumé et les 
légendes des figures.

L'inclusion d'une image simple est montrée sur la figure \ref{logo},
et un agencement plus sophistiqué sur les figures \ref{gmt} à \ref{draftimage}.

\begin{figure}
\centerline{\includegraphics[width=0.4\linewidth]{LogoIPGPUPC}}
\caption{Logo de l'IPGP depuis juin 2022.}
\label{logo}
\end{figure}


\begin{minipage}{\linewidth}
Deux erreurs sont très communes 
\begin{enumerate}
\item Des figures inutilement grandes : la résolution d'une figure 
est plus importante que la taille et souvent une même figure est 
plus ``sympathique'' en plus petit.
\item Des figures tournées de 90$^\circ$ qui forcent le lecteur à 
la tourner, pour l'afficher en mode ``paysage''. 
C'est très rare qu'il existe des bonnes raisons à cela.
\end{enumerate}
\end{minipage}

\begin{figure}

\begin{minipage}[b]{0.48\linewidth}
\includegraphics[width=\linewidth]{gmt1}
\caption{Image produite par \href{https://www.generic-mapping-tools.org/}{GMT}.}
\label{gmt}
\end{minipage}
\hfill
\begin{minipage}[b]{0.48\linewidth}
\fboxsep 0pt
\centerline{\fbox{\includegraphics[bb=177 146 200 161,clip=,width=0.75\linewidth]{gmt1}}}
\caption{Zoom sur la Fig~\ref{gmt}.}
\label{zoomimage}
\end{minipage}

\vspace*{5mm}

\begin{minipage}{0.48\linewidth}
\includegraphics[width=\linewidth,draft]{gmt1}
\end{minipage}
\hfill
\begin{minipage}[b]{0.48\linewidth}
\caption{Comme la Fig~\ref{gmt}, mais en \texttt{draft}.}
\label{draftimage}
\end{minipage}

\end{figure}

\subsection{Les tables}

L’environnement pour créer des tableaux est \texttt{tabular}. Cet environnement
doit être accompagné d'un argument qui précise
le nombre de colonnes, et la position de leur contenu (à droite, centrée ou à gauche).
Le passage d'une colonne à l'autre se  fait avec le symbole \& et le passage
d'une ligne à l'autre par \verb"\\". Il est possible d'ajouter des lignes verticales.
(dans la définition des colonnes) et des lignes horizontales avec \verb"\hline" après le \verb"\\".

Les commandes \verb"\toprule", \verb"\midrule" et \verb"\bottomrule" du package \verb"booktabs"
permettent d'obtenir des lignes plus esthétiques que les lignes standard.


Un tableau peut être encapsulé dans un environnement \texttt{table}, qui est un objet flottant comme les figures,
mais il est d'usage de mettre la légende avant le tableau lui-même.

\begin{table}
\caption{Quelques nombres  de pétales.}
\centerline{
\begin{tabular}{lr}
\toprule
\textbf{Fleur} & \textbf{Pétales (ou fleurons)} \\
\midrule
Marguerite &    21 \\
Tulipe & 6 \\
Rose sauvage & 5 \\
Lys & 6 \\
\bottomrule
\end{tabular}}
\end{table}


\subsection{Les équations}

Écrivons quelques lignes de texte technique pour illustrer la façon 
normale d'écrire des équations.

Soit un milieu parfaitement élastique, et soient
\ $ \sigma^{ij} ({\bf x},t) $ \ 
et
\ $ \varepsilon^{ij} ({\bf x},t) $ \ 
respectivement les composantes 
\ $ \{i,j\} $ \ 
de la contrainte et de la déformation au point 
\ $ {\bf x} $ \ 
et à l'instant 
\ $ t $~.
Pour des petites déformations, nous pouvons écrire, au premier ordre,
\begin{equation}
     \sigma^{ij} ({\bf x},t) 
     \ = \  
     \sigma_0^{ij} ({\bf x},t) 
     + 
     c_0^{ijkl} ({\bf x}) 
     \,  
     \varepsilon^{kl} ({\bf x},t) 
     \ , 
          \label{eq:loiHooke}
\end{equation}
où le terme 
\ $ \sigma_0^{ij} ({\bf x},t) $ \ 
est appelé \emph{précontrainte\/}, et où les coefficients 
\ $ c_0^{ijkl} ({\bf x}) $ \ 
sont les \emph{constantes élastiques\/}. 
L'équation~(\ref{eq:loiHooke}) représente la \emph{ Loi de Hooke\/}.
Maintenant, que représente l'équation
\begin{equation}
     E({\bf x},t) 
     \ = \  
     \frac{1}{2} 
     \, 
     \sigma^{ij} ({\bf x},t) 
     \, 
     \varepsilon^{ij} ({\bf x},t)
     \ ? 
          \label{eq:energydensity}
\end{equation}
Peu importe en vérité, car elle est hors sujet. 
Par contre, en utilisant la loi de Hooke~(\ref{eq:loiHooke}) et le principe
 fondamental de la dynamique, 
on démontre l'\emph{équation des ondes élastiques\/} :
\begin{equation}
     \rho \ 
     \frac{\partial^2 u^i}{\partial t^2} 
     - 
     \frac{\partial}{\partial x^j} 
     \left( 
           c_0^{ijkl} 
           \, 
           \frac{\partial u^k}{\partial x^l} 
     \right)
     \ = \  
     0 
     \ . 
          \label{eq:ondes}
\end{equation}

La discussion ci-dessus illustre quelques unes des règles à utiliser dans 
un texte contenant des équations :

\begin{itemize}
     \item  Il faut introduire dans le texte {\bfseries toutes} les variables.

     \item  Il faut numéroter {\bfseries toutes} les équations.

     \item  Dans la mesure du possible, il faut donne un \emph{nom\/} à 
            chaque équation, et, dans la suite du texte, dire 
            ``l'équation des ondes~(\ref{eq:ondes})'' 
            au lieu de ``l'équation~(\ref{eq:ondes})''.

     \item  Il faut \emph{ponctuer\/} les équations, car elles font 
            partie du texte. Ici, non seulement nous avons des équations 
            se terminant par des virgules et des points, mais même 
            une équation se terminant par un point d'interrogation.

     \item  N'abusez pas du ``~:~'' pour arriver aux équations : 
            dans les exemples ci-dessus, seulement une équation le permet.
\end{itemize}

\subsection{La typographie}

Soignez l'orthographe. Évitez même les erreurs comme celle de ne pas accentuer 
les majuscules (cette phrase commence par une majuscule accentuée, et pour vous
 con\-vain\-cre, regardez les couvertures de livres ou les titres des
journaux, ou vérifiez dans votre dictionnaire que l'on écrit bien 
États-Unis\dots).

\subsection{\LaTeX, Word, ou quoi d'autre?}

Tout logiciel de préparation de documents peut être utilisé pour produire 
votre rapport. \LaTeX \ est bien adapté à la publication scientifique, 
avec ses notations ma\-thé\-ma\-ti\-ques. Mais, surtout, il est devenu le 
standard pour la publication technique.

Pour votre rapport vous devrez utiliser exactement le style illustré par ce document. 
Si vous utilisez \LaTeX, vous n'aurez rien de spécial à faire : 
ni à définir les 
 polices de caractères (ou fontes), ni les espaces entre sections, etc.), car 
ceci n'est qu'une adaptation mineure du style ``article'' de \LaTeX. 
Je déconseille de choisir un autre logiciel, moins standard ou moins technique.
Si vous \^{e}tes obligés de le faire, reproduisez exactement le style de ce 
document. Pensez à bien utiliser une taille de police de caractères de 11~points 
\footnote{Au passage, pour éviter de couper entre un nombre et son unité, 
utiliser le tilde (blanc insécable) : \texttt{11$\sim$points}.} 
pour le texte principal. 

\subsection{Les erreurs courantes  à éviter}

Une erreur courante des auteurs inexpérimentés est de ne pas donner assez
 d'importance aux références bibliographiques : références manquantes, 
références citées dans le texte, mais absentes de la liste finale, etc. 
Souvent aussi, la liste des références ne suit pas un des formats habituels, 
ou il y manque un des éléments (année, page, etc.).

Souvent les figures sont trop grandes, ou bien elles dépassent des marges de 
l'article. 
Évitez d'obliger le lecteur à tourner la revue : il est rare qu'une figure
 doive absolument avoir un format paysage sur une page entière.
Parfois il y a une multiplicité de légendes : on écrit des phrases dans la
 figure elle-même.      

Il arrive que le résumé ou l'introduction manquent, ou que l'on confonde 
leurs rôles respectifs.     

Si un document a un seul auteur, on ne doit jamais utiliser le ``nous'',  
qui est réservé au Pape. Le ``je'' doit être utilisé dans des très 
rares occasions. Soyez neutre : ``il'' est possible de montrer que\dots 

Une grave erreur consiste à sous-estimer le temps qu'il faut pour
 produire un rapport sous sa forme finale : commencez à rédiger dès la mi-stage.
Même si vos recherches sont loin d'avoir abouti, vous pouvez déjà commencer
 à mettre la bibliographie ---\,et, surtout, vos idées\,--- en place.
 
 
 Autres erreurs communes :
\begin{itemize}
\item utiliser le retour chariot (\verb"\\") pour finir un paragraphe,
\item essayer de forcer le placement des objets flottants comme les figures et les tables (option \verb"[h]"),
\item continuer à saisir du texte, sans corriger les erreurs de compilations, et ne pas regarder les causes d'erreurs dans les fichiers log.
\end{itemize}



\subsection{Utilisation de \textsc{Bib}\LaTeX{} avec (\texttt{biber}) pour les références}
Il est impératif que vous preniez l'habitude d'utiliser \textsc{Bib}\LaTeX{} 
pour gérer votre
 base de données de références et son adressage dans votre texte. L'utilisation de \textsc{Bib}\LaTeX{} 
est 
 détaillée dans le texte de \cite{moguilnyIntro}. Ne pas oublier d'ajouter les DOI 
 (\textit{Digital Object Identifier}) dans les entrées de la base de données \texttt{.bib}.

 Quelques emplois possibles :

\verb"\textcite{robert2012theory}"  donnera \textcite{roberts2012theory}  \\
\verb"\parencite{roberts2012theory}"  donnera \parencite{roberts2012theory}  \\
\verb"\parencite[voir par exemple][\S 2]{roberts2012theory}"  donnera \\
\hspace*{1cm}\parencite[voir par exemple][\S 2]{roberts2012theory}  \\
\verb"\citeauthor{roberts2012theory}"  donnera \citeauthor{roberts2012theory}  \\
\verb"\citeyear{roberts2012theory}"  donnera \citeyear{roberts2012theory}  

La variable \texttt{pubstate} dans l'entrée d'une publication de la base de données permet de préciser son état, par exemple
``soumis'' pour \textcite{frezatsoumis}. 

\section{Conclusion}

Écrire un rapport de stage de master n'est pas aisé. 
Cela prend du temps, et les difficultés sont souvent sous-estimées.
Mais cette rédaction est une composante essentielle du stage de recherche.

Le \emph{style\/} d'un rapport technique est bien codifié, et laisse peu de place
à l'originalité, qui, par contre, aura libre cours dans le choix du sujet
ou des méthodes de recherche. On demande une grande \emph{ clarté\/} et une grande 
\emph{concision\/} dans la présentation. 
Pour cela, il faut s'y prendre avec du temps : un document écrit se travaille 
beaucoup avant de le divulguer.

\section{Remerciements}

Généralement, vous serez le seul signataire du rapport ---\,et, donc, le seul 
responsable\,---.
Votre directeur de stage n'apparaîtra que dans cette section 
``\emph{Remerciements\/}''. 
Si votre rapport contient des résultats scientifiques, veillez à faire justice, 
et faire apparaître ici les apports respectifs de votre directeur de stage
et de vous même.

Un style typique pour cette section est le suivant :

\paragraph*{Remerciements.} Ce stage de master a été effectué dans le 
Laboratoire de Géo\-pan\-tou\-flis\-me de 
l'Institut de physique du globe de Paris, sous la direction de Napoléon Bonaparte, 
que je voudrais remercier pour la qualité du café qu'il sait faire. 
Émile Zola m'a aidé avec le logiciel \LaTeX \ de préparation de documents. 
Le logiciel de dessin utilisé appartient à l'Université de Cochabamba.
Cette recherche a été menée à bien grâce à une bourse du 
Ministère des Anciens Ministres, et le microscope aquatique utilisé a 
été financé par le Centre National de la Recherche Parapsychologique. 
 
\newpage

\addcontentsline{toc}{section}{Références}
%\nocite{*} % pour afficher toutes les entrées (même celles non citées)
            % de la base de données bibliographique

\printbibliography
            
\newpage

\appendix % (Annexes) si nécessaire / if necessary
\section{Annexes}
 
\subsection{Précisions sur un certain point}
  
\subsection{Précisions sur un autre point}
 
\cleardoublepage
  
\end{document}

